% Proved after four complete actual-file and exact-source audits. Arbitrary covers and arbitrary field characters.
\subsection{Hereditary semilinear growth with arbitrary field monodromy}

\begin{lemma}[Proved: equivariant comparable cycles on the field cover]
\label{D:equivariant-cover-comparable-cycles}
Retain the fixed complete normal surface, base resolution $Y_0$,
parameter cycle $C_0$, and constants of
Proposition~\ref{D:uniform-comparable-reduction-cycle}.
Let $U_V\to U_R$ be a connected finite etale cover of degree $s$,
and let a finite group $G$ act on this cover over $U_R$.
The regular model used in that proposition can be chosen with
the lifted $G$-action. On this same model the actual cycles
\[
 C=g^*C_0,\qquad Z=\alpha C-\beta R_g,\qquad
                 C\le Z\le\alpha C                         \tag{D.SH.1}
\]
are $G$-invariant. All scalar, generation, and canonical
Frobenius bounds of that proposition and
Theorem~\ref{D:arbitrary-cover-ordinary-hereditary-growth}
hold with their original constants and cover degree $s$.
No invertibility assumption on $|G|$ is needed.
\end{lemma}

\begin{proof}
Let $X$ be the finite normalization of $Y_0$ in the cover's
function field. The action extends to $X$ by the integral
closure property. The model in the comparable-cycle theorem
is the minimal regular resolution of $X$: a curve maps to a
point of $Y_0$ exactly when it is exceptional over $X$, since
$X\to Y_0$ is finite. Thus contracting precisely the
point-mapped exceptional curves of the first kind gives
the relative-minimal model stipulated in that theorem.

We spell out the equivariance deduction from the standard
surface facts, rather than require an equivariant resolution
algorithm. Resolve a rational map between two minimal regular
resolutions of $X$ by a shortest sequence of point blowups.
The other projection is also a sequence of point blowups.
These are the known statements of
\citet[Lemmas 54.4.3 and 54.17.1, Tags 0C5H and 0C5R]
{StacksSkewMinimalModelD}. Let $E$ be the last exceptional
$(-1)$-curve of the first projection. If the other projection
contracts $E$, its factorization through that contraction
contradicts shortest length, by
\citet[Lemma 54.16.1, Tag 0C5J]{StacksSkewMinimalModelD}.
Otherwise its image is a curve exceptional over $X$.
Successive blowdowns increase its square from $-1$ by sums
of nonnegative squares. Negative definiteness of the
exceptional intersection form makes the final square
negative, so every increase is zero. Hence no blowdown
centre meets the curve; its final image is still a smooth
rational $(-1)$-curve. This contradicts minimality.
There were therefore no blowups, and interchanging the
resolutions shows that the map is an isomorphism.
An isomorphism over the normal base is unique because it is
the identity on the common dense regular locus.

Every automorphism of $X$ consequently lifts uniquely; the
group identities follow from uniqueness. This argument is
independent of the characteristic and of the group order.
The map $g:Y\to Y_0$ is invariant. Its actual effective
Jacobian divisor $R_g$ is invariant by functoriality of the
map of canonical sheaves. Thus both cycles in (D.SH.1) are
invariant. This is exactly the same relative-minimal model
on which the bound $R_g\le\lambda C$ and the all-complete
degree-linear geometric budgets were proved. No additional
equivariant blowups, and no enlargement of their constants,
have been introduced.
\end{proof}

\begin{theorem}[Proved: arbitrary-cover hereditary semilinear growth]
\label{D:arbitrary-cover-semilinear-hereditary-growth}
Let $(R,\mathfrak m,k)$ be a complete excellent normal local
surface with algebraically closed coefficient and residue
field $k$ of characteristic $p>0$. Put
$\Gamma=\pi_1(\operatorname{Spec}R\setminus\{\mathfrak m\})$.
There is a constant $A_R^{\rm sl}$ such that, for every open
subgroup $H\subset\Gamma$ of index $t$, every $q=p^f$,
every $r\ge1$, and every fixed continuous character
\[
 \chi:H\longrightarrow
                 \operatorname{Gal}(\mathbb F_q/\mathbb F_p),
                                                               \tag{D.SH.2}
\]
the representations
\[
 H\longrightarrow \operatorname{GL}_r(\mathbb F_q)
          \rtimes\operatorname{Gal}(\mathbb F_q/\mathbb F_p)
 \quad\text{with field projection }\chi
                                                               \tag{D.SH.3}
\]
have at most $q^{A_R^{\rm sl}tr^2}$ conjugacy classes under
$\operatorname{GL}_r(\mathbb F_q)$.
There is no tameness condition on either cover, no bound
on the field-image order, and no assumption that the
field character extends over a resolution.
\end{theorem}

\begin{proof}
Write $d=|\operatorname{im}\chi|$, $\nu=f/d$, and
$q_0=p^\nu$. The subgroup $K=\ker\chi$ gives a connected
cyclic degree-$d$ cover $U_V\to U_H$. Its total degree over
$U_R$ is $td$. The finite normalization $A_V$ is complete,
normal, and local with residue field $k$. Choose the deck
generator $\tau$ with the pullback convention of
Lemma~\ref{D:twisted-deck-Frobenius-semilinear-model}.
By that proved category equivalence, (D.SH.3) is exactly
the category of bundles $E$ on $U_V$ with
\[
              \Phi:\tau^*F_{q_0}^*E\xrightarrow{\sim}E.
                                                               \tag{D.SH.4}
\]
The equivalence includes morphisms and arbitrary coefficient
base extension. The $d$-fold composite defines the ordinary
$F_q$-unit structure; it is not an independently imposed
matrix equation or a false order-$d$ relation on a chosen
element of $H$.

Use Lemma~\ref{D:equivariant-cover-comparable-cycles} for
the cyclic deck group on the degree-$td$ cover. We therefore
have a deck-equivariant model $Y$ and invariant $C,Z$ with
all the ordinary theorem's budgets $O_R(td)$.
In particular, $C\le Z\le\alpha C$,
$\mathcal O_Y(-C)$ and $\mathcal O_Y(-Z)$ are generated,
and $H^1(Z,\mathcal O_Z(-jZ))=0$ for every $j\ge1$.

\emph{Short Frobenius cycles and determinant choices.}
Factor (D.SH.4) into $\nu$ one-step arrows: for
$0\le a<\nu-1$ they are $F_p^*E_a\to E_{a+1}$,
and the last arrow is $\tau^*F_p^*E_{\nu-1}\to E_0$.
The intermediate punctured bundles can be $F_p^{a*}E$.
For geometric input representations take their reflexive
extensions and canonical full pullbacks $F_a$ on $Y$.
The forward arrows are regular and the inverse arrows
have poles at most $pC$, exactly as in the ordinary proof.
This includes the last arrow: canonical pullback commutes
with the automorphism $\tau$, and $\tau^*C=C$.

The determinant of (D.SH.3) is a crossed character with
coefficients $\mathbb F_q^\times(\chi)$. By
Lemma~\ref{D:arbitrary-field-character-scalar-H1} there
are at most $q^{c_Rt}$ such classes. This enumeration is
made on $H$, before passing to the field cover. Fix one
class; its restriction to $U_V$ fixes the ordinary
punctured determinant of $E$, and its $p^a$-powers fix
the determinants at the other $\nu$ stages.

Put $x_a=c_1(F_a)\cdot C$. These numbers are nonnegative.
The determinant divisors of the forward arrows lie
between zero and $rpC$. At the last stage
$c_1(\tau^*F_{\nu-1})\cdot C=x_{\nu-1}$.
The same cyclic recurrence as in the ordinary theorem
therefore gives
\[
     0\le x_a\le\frac{rp}{p-1}(-C^2)=O_R(td r).
                                                               \tag{D.SH.5}
\]
Every coefficient of $C$ is a positive integer. Hence
the unweighted exceptional determinant multidegrees have
nonnegative entries with sum $O_R(td r)$.
There are $O_R(td)$ components, and the negative definite
intersection matrix makes the multidegrees determine
the extension of a fixed punctured line uniquely.
The binomial enumeration in the ordinary theorem bounds
the ordered lists of global lines
$L_a=\det(F_a(-Z))$ by
\[
 q^{c_Rt}\exp\bigl(c_1\nu td\log(r+1)\bigr)
       =q^{c_Rt}\exp\bigl(c_1ft\log(r+1)\bigr).         \tag{D.SH.6}
\]
No list of all ordinary characters on the degree-$td$
cover is substituted for the twisted scalar count.

\emph{The actual finite polynomial parameter scheme.}
Choose $N\ge\max\{3,\alpha\}$, $L=N+p$, and $D=LC$,
as in the ordinary proof. All these Cartier subschemes
are deck-invariant. For a fixed ordered line tuple use
the generated quotient charts
\[
 0\longrightarrow L_a^{-1}\longrightarrow
 \mathcal O_D^{r+1}\longrightarrow V_a\longrightarrow0,
             \qquad E_a=V_a(Z).                        \tag{D.SH.7}
\]
They are actual locally split sequences over every
coefficient algebra. The scalar cohomology and line-degree
budgets on $D$ are $O_R(td)$ and $O_R(td r)$.
Their object coordinates are $O_R(td r^2)$ per stage.

For the first $\nu-1$ arrows impose the two inverse
relations of the ordinary theorem. For the last arrow
put $P=\tau^*F_p^*E_{\nu-1}$ and introduce maps
\[
 a:P\longrightarrow E_0,\qquad
 b:E_0\longrightarrow P(pC),\qquad
 a(pC)b=\iota_{E_0},\quad ba=\iota_P.                 \tag{D.SH.8}
\]
Both equalities are equalities of maps to the indicated
Cartier twists on $D$; the Cartier equation is not inverted
on a nonreduced finite cycle.

Here is the precise effect of the deck pullback on the
ordinary coordinate proof. It replaces the fixed line
$L_{\nu-1}$ by $\tau^*L_{\nu-1}$ in the last source.
This line is generated and has the same degrees as the
original line. It fixes the geometric lines
$\mathcal O_D(-C)$ and $\mathcal O_D(-Z)$.
The pulled-back quotient tuple has the same parameter
coefficients; its only nonlinear operation is the one
$p$-th power from Frobenius. All reference sections and
their deck pullbacks are fixed background data.
One may compute all required Hom complexes on any fixed
two-affine cover of $D$, treating the pulled-back fixed
lines as additional reference lines. No deck-equivariant
choice of cohomology bases is required.

The four-term quotient-Hom calculation of
Proposition~\ref{D:fixed-determinant-finite-Hom}, with the
bounded powers of generated lines used in the ordinary
proof, still has cohomology dimension $O_R(td r^2)$.
Its exact all-ring cycle, composition, and boundary-witness
equations have the same bounded polynomial degrees.
There are a fixed number of complexes per stage, including
(D.SH.8). Thus the whole affine zero scheme has
\[
 \#\text{coordinates},\ \#\text{scalar equations}
       \le c_2\nu td r^2=c_2ft r^2,
 \qquad \deg(\text{each equation})\le D_R,             \tag{D.SH.9}
\]
with constants independent of $d,t,r,f$.
Its generating-tuple conditions form one open subset.
The fixed reference cochains may depend on the cover and
the line tuple, but their coefficients are constants and
do not change variable degrees or coordinate counts.

\emph{Effectivity of the skew cycle over coefficient bases.}
The local matrix construction of
Proposition~\ref{D:finite-Frobenius-object-effectivity}
extends verbatim at the level of its specified formulas.
For precision, take a finite affine trivializing cover
closed under $\tau$; the number of this qualitative local
cover is not a parameter count. At the last step the
correction operator on transition matrices is
$g\mapsto A_i\,\tau^*F_p^*(g)\,A_j^{-1}$, with the
appropriate permuted affine indices. The inverse still
has pole at most $pC$. Since $\tau$ preserves the ideal
of $C$, an error of order $j$ becomes an error of order
at least $pj-p$. Thus $pN-p>N$ gives the same convergent
correction, unique with its prescribed $NC$ marking.
Products and inverse transitions satisfy the same
correction and marking, so uniqueness proves the actual
cocycle relations. This works over every coefficient
algebra, with no averaging and no invertibility of $d$.
The resulting formal bundles have the actual skew unit
arrows. On finite-type coefficient bases they algebraize
over the proper Noetherian completed tensor base by the
same Grothendieck-existence argument. Arbitrary-ring
finite-cycle and formal assertions remain separate from
that Noetherian algebraization assertion.

The retained $NC$ marking includes $Z$. It therefore
retains the generating quotients and the determinant
lines in (D.SH.7). Lemma
\ref{D:weak-cycle-generated-determinant-lifting} lifts
these generators and identifies the formal determinant
with the fixed global line. This argument does not depend
on the Frobenius cycle length or on its final deck pullback.
It makes no relative canonicality or coflatness assertion.

\emph{The actual punctured class comparison.}
The proof of
Proposition~\ref{D:generated-fixed-determinant-unit-openness}
applies to these generated, fixed-determinant families.
Indeed target Hartogs and adjugates bound both directions
of an Artinian punctured isomorphism by a fixed multiple
of $C$, independently of its nilpotence order.
The multiple may depend on the reference class and rank;
it is used only for qualitative finite presentation.
The intrinsic finite unit-Hom equation has the same
order improvement at the final arrow because $\tau^*$
preserves $C$. The shifted Hom and product lattices,
paired inverse equations, and uniqueness of their formal
corrections are consequently unchanged.

For the infinitesimal lifting assertion, iterate the skew
cycle $d$ times only to identify its finite etale frame
sheaf with the ordinary $F_q$-unit frame sheaf. The
semilinear deck descent is the one in the proved category
equivalence (D.SH.4). A prescribed unit isomorphism over
a nilpotent quotient lifts uniquely by finite etale
invariance; its deck compatibility follows by the same
uniqueness. The preceding generated-family argument puts
this lift inside the fixed meromorphic bound. Hence the
represented comparison Isom functor is locally of finite
presentation and formally etale, so its class images are
open and closed, exactly as in the ordinary theorem.
The $d$-fold composite is not expanded in the polynomial
equations (D.SH.9).

Every original representation is covered. Its canonical
input objects give (D.SH.7)--(D.SH.8) for one determinant
tuple, and its genuine formal transition matrices are
one admissible choice in the correction construction.
Choice-independence identifies the constructed object
with that original punctured skew unit class.

Finally the rectangular affine component bound, applied
to (D.SH.9), gives at most
$(M+1)\max\{2,D_R\}^M\le p^{c_3ft r^2}$ classes for
each line tuple, where $M\le c_2ft r^2$.
Multiplying by (D.SH.6), and using
$\log(r+1)=O(r^2)$ and $\log q=f\log p$, proves
$q^{A_R^{\rm sl}tr^2}$.
The category equivalence identifies these classes with
exactly the conjugacies in (D.SH.3). This also treats
$d=1$ and allows $p\mid d$.
\end{proof}

\begin{remark}[Scope of the field extension]
The theorem removes the resolution-extension and radial-tame
conditions on the field character in the earlier semilinear
counts. It uses the newly proved twisted scalar bound rather
than an all-characteristics local $H^1$ theorem. Its target
is the ordinary general semilinear group with a prescribed
field projection. No projective lifting or twisted-gerbe
representation count is inferred.
\end{remark}
